% TEMPLATE FOR A NEW CONCEPT PAGE
%
% 1. Copy this file to sections/foundations/<id>.tex, where <id> is a short name (e.g. maximum-principle).
%    The same name is the page's address on the website: notes/pde/<id>.html.
% 2. Replace every <...> below.
% 3. In main.tex, add  \input{sections/foundations/<id>}  under the chapter (layer) it belongs to.
%    The order of the \input lines is the reading order (previous / next links on the website).
%
% Links to other pages: \pref{gronwall}{Gronwall's inequality}  prints "Gronwall's inequality (§2.4)"
% in the PDF and is a link on the website. \pref{gronwall}{} uses the page's own title.

\section{<Title>}\label{p:<id>}

\status{stub}                      % stub (statement + reference only), draft, or done
\prereqs{<id-1>, <id-2>}           % pages to read first; delete the line if there are none

\subsection*{The idea}

<What the tool is for, in plain words.>

\subsection*{Statement}

\begin{thmbox}{Theorem (<name>)}
<The precise statement, with every hypothesis.>
\end{thmbox}

\subsection*{Why it is plausible}

<A heuristic instead of a proof: a scaling argument, a picture, a model example.>

\subsection*{Exercises}

\begin{exercise}
  <A problem that makes the reader do the idea by hand.>
  \begin{hint}
    <The key step, without the full solution.>
  \end{hint}
  \begin{solution}
    <The solution.>
  \end{solution}
\end{exercise}

\subsection*{Where to find a proof}

\begin{itemize}
\item <Author>, \emph{<Book>}, <section and theorem number>.
\end{itemize}

\neededfor                          % "Needed for": filled in from the other pages' \prereqs
